\documentclass[11pt,a4paper]{article}
\usepackage[utf8]{inputenc}
\usepackage[T1]{fontenc}
\usepackage[spanish,es-nodecimaldot]{babel}
\usepackage{lmodern}
\usepackage{geometry,booktabs,tabularx,array,enumitem,amsmath,graphicx,float,caption,hyperref}
\geometry{top=2.3cm,bottom=2.3cm,left=2.4cm,right=2.4cm}
\hypersetup{colorlinks=true,linkcolor=black,urlcolor=blue}
\setlength{\parindent}{0pt}
\setlength{\parskip}{0.55em}
\setlength{\tabcolsep}{4pt}
\renewcommand{\arraystretch}{1.15}
\setlist{leftmargin=*,topsep=3pt,itemsep=3pt}
\newcolumntype{Y}{>{\raggedright\arraybackslash}X}
\emergencystretch=2em
\raggedbottom

\title{\textbf{Exploración de modelos de inteligencia artificial multimodal\\
para la identificación de macromicetos mediante dispositivos móviles}\\[0.6em]
\normalsize Resultados de la semana 11: comparación de variantes y selección del modelo}
\author{Daniel Esteban Camacho Ospina\\Cesar Augusto Pulido Cuervo}
\date{}

\begin{document}
\maketitle

\begin{abstract}
Se compararon seis variantes para clasificar diez géneros de macromicetos y se
evaluó su sensibilidad a datos contextuales ausentes, imprecisos o incorrectos.
Imagen, GPS y clima obtuvo el mayor F1 macro en validación, con 0,971950.
La diferencia frente a solo imagen fue pequeña y su intervalo de confianza
incluyó cero. Se seleccionó la variante visual como configuración operativa.
Se presentan curvas del entrenamiento visual, métricas por género, matrices
de confusión y ejemplos Grad-CAM de aciertos y errores.
\end{abstract}

\section{Objetivo y protocolo}
Todas las variantes comparten MobileNetV3-Large S5-C, congelada. Se ajustaron
cinco regresiones logísticas balanceadas con 17\,799 observaciones de
entrenamiento de FungiTastic \texttt{v002}. Se promediaron las probabilidades
visuales de las fotos por observación. Cuando hay contexto utilizable, la
fusión combina 90\,\% de probabilidad visual y 10\,\% de probabilidad contextual.
La validación conserva las mismas 3\,452 observaciones para las seis variantes.
A--D reproducen las semanas 9 y 10. E permite estudiar el clima sin GPS y F
permite estudiar el sustrato sin clima. Cada rama contextual se ajustó con las
variables de su combinación.

\section{Comparación en validación}
\begin{table}[H]
\centering
\caption{Resultados por observación en validación.}
\label{tab:val}
\begin{tabularx}{\linewidth}{lYrrrr}
\toprule
Modelo & Modalidades & F1 macro & Exactitud (\%) & Aciertos & Errores \\
\midrule
A & Imagen & 0,971674 & 97,68 & 3\,372 & 80 \\
B & Imagen + GPS & 0,971674 & 97,68 & 3\,372 & 80 \\
C & Imagen + GPS + clima & 0,971950 & 97,71 & 3\,373 & 79 \\
D & Imagen + GPS + clima + sustrato & 0,971880 & 97,71 & 3\,373 & 79 \\
E & Imagen + clima & 0,971848 & 97,68 & 3\,372 & 80 \\
F & Imagen + GPS + sustrato & 0,971674 & 97,68 & 3\,372 & 80 \\
\bottomrule
\end{tabularx}
\end{table}

\clearpage
El F1 macro promedia los F1 de los diez géneros con igual peso. C corrigió dos
errores de A e introdujo uno. La diferencia de F1 macro fue de
+0,028 puntos en la escala 0--100, con un intervalo de confianza del
95\,\% de -0,036 a +0,101. El intervalo se calculó con 2\,000 remuestreos pareados
por observación.

\section{Sensibilidad a datos faltantes e imprecisos}
\begin{table}[H]
\centering
\caption{Cambio de F1 macro en puntos sobre escala 0--100 respecto a la condición original.}
\label{tab:sens}
\begin{tabular}{lrrrrr}
\toprule
Condición & B & C & D & E & F \\
\midrule
Todo el contexto ausente & +0,000 & -0,028 & -0,021 & -0,017 & +0,000 \\
GPS ausente & +0,000 & +0,117 & +0,084 & --- & -0,085 \\
GPS redondeado a 1 grado & +0,000 & +0,059 & +0,000 & --- & +0,000 \\
GPS permutado & +0,000 & -0,012 & -0,012 & --- & +0,006 \\
Clima ausente & --- & +0,017 & +0,108 & -0,017 & --- \\
Clima permutado & --- & -0,010 & -0,014 & +0,011 & --- \\
Sustrato ausente & --- & --- & -0,001 & --- & +0,018 \\
Sustrato permutado & --- & --- & +0,000 & --- & +0,000 \\
\bottomrule
\end{tabular}
\par\smallskip{\footnotesize El guion indica que la modalidad queda fuera de la variante.}
\end{table}

La ausencia total activa la rama visual y reproduce A. C gana 0,117 puntos al
retirar el GPS y D gana 0,108 al retirar el clima en esta cohorte. El redondeo
a un grado modifica el GPS y conserva el clima asociado. La permutación
reasigna bloques de metadatos entre observaciones con semilla 42 y simula
metadatos incorrectos.

\section{Evaluación externa y selección}
\begin{table}[H]
\centering
\caption{Transferencia externa por observación.}
\label{tab:ext}
\begin{tabular}{lrrrr}
\toprule
Modelo & F1 total & Exactitud total (\%) & F1 Colombia & F1 Orinoquía \\
\midrule
A & 0,766287 & 85,78 & 0,772384 & 0,576427 \\
B & 0,762022 & 85,25 & 0,763822 & 0,578913 \\
C & 0,762022 & 85,25 & 0,763822 & 0,578913 \\
D & 0,765217 & 85,72 & 0,769714 & 0,576427 \\
E & 0,766287 & 85,78 & 0,772384 & 0,576427 \\
F & 0,762063 & 85,25 & 0,763848 & 0,579860 \\
\bottomrule
\end{tabular}
\par\smallskip{\footnotesize El F1 macro promedia los diez géneros en todos los grupos.}
\end{table}

El conjunto externo contiene 4\,895 observaciones ya evaluadas en S9--S10.
Colombia (3\,348) incluye Orinoquía (309). Existe GPS para parte de los registros.
Sin clima ni sustrato compatibles, E reproduce A y el aporte de las modalidades
completas en Colombia queda por medir. El test interno quedó fuera de este estudio.

Se selecciona \textbf{A, solo imagen}, como configuración operativa por su menor
dependencia de metadatos y la incertidumbre de la pequeña ganancia interna de C.
El conjunto externo se usa como diagnóstico de transferencia.

\section{Evaluación del modelo visual}
\subsection{Durante el entrenamiento}
\begin{figure}[H]
\centering
\includegraphics[width=\linewidth]{entrenamiento.png}
\caption{Pérdida de entrenamiento y exactitud de validación por fotografía.
El marcador identifica la época 16, elegida por F1 macro de validación por
observación.}
\label{fig:curvas}
\end{figure}

En 21 épocas, la pérdida de entrenamiento bajó de 1,343 a 0,652 y la
exactitud de validación por fotografía subió de 89,32\,\% a 95,02\,\%.
En la época seleccionada, la exactitud por foto fue de 94,29\,\%. A partir
de las probabilidades archivadas se estimó una pérdida de validación de
0,829 para esa época, usando entropía cruzada ponderada y suavizado de
etiquetas. La estimación agrega la validación completa y usa probabilidades
redondeadas. El historial carece de pérdida de validación y exactitud de
entrenamiento por época. El diagnóstico de sobreajuste mediante la divergencia
entre pérdidas requiere ambas curvas.

\subsection{Después del entrenamiento}
La precisión mide qué proporción de las predicciones de un género es correcta.
La sensibilidad (\emph{recall}) mide qué proporción de las observaciones reales
de ese género detecta el modelo. El F1 es la media armónica de ambas.
\begin{table}[H]
\centering
\caption{Precisión, sensibilidad y F1 de A en el conjunto externo completo.}
\label{tab:genera}
\begin{tabular}{lrrrr}
\toprule
Género & Precisión (\%) & Sensibilidad (\%) & F1 & Soporte \\
\midrule
Amanita & 95,97 & 98,45 & 0,97196 & 1\,549 \\
Auricularia & 85,47 & 50,87 & 0,63774 & 578 \\
Ganoderma & 80,59 & 88,89 & 0,84536 & 369 \\
Geastrum & 92,49 & 93,81 & 0,93144 & 420 \\
Hericium & 13,64 & 100,00 & 0,24000 & 6 \\
Lycoperdon & 56,74 & 77,69 & 0,65584 & 130 \\
Morchella & 96,57 & 96,57 & 0,96571 & 175 \\
Pleurotus & 74,33 & 84,34 & 0,79022 & 364 \\
Trametes & 77,74 & 90,26 & 0,83536 & 801 \\
Xylaria & 90,49 & 69,98 & 0,78924 & 503 \\
\bottomrule
\end{tabular}
\end{table}

\clearpage
\subsection{Matrices de confusión}
Las matrices de las figuras~\ref{fig:val} y~\ref{fig:ext} muestran los aciertos
en la diagonal y los errores fuera de ella. En el externo, las confusiones más
frecuentes de A fueron Auricularia como Trametes (106 observaciones) y como
Pleurotus (76).

\begin{figure}[H]
\centering
\includegraphics[width=0.94\linewidth]{confusion.png}
\caption{Matrices de confusión de las seis variantes en validación por observación.
El color codifica el recuento de errores.}
\label{fig:val}
\end{figure}

\begin{figure}[H]
\centering
\includegraphics[width=0.94\linewidth]{confusion_externo.png}
\caption{Matrices de confusión de las seis variantes en el conjunto externo.
La escala de color se satura a partir de 12 errores.}
\label{fig:ext}
\end{figure}

\subsection{Grad-CAM}
\begin{figure}[H]
\centering
\includegraphics[width=0.95\linewidth]{gradcam.png}
\caption{Recorte de entrada y Grad-CAM para el género predicho en tres
fotografías externas de iNaturalist. La selección por identificador incluye un
acierto y dos errores de géneros reales distintos.}
\label{fig:gradcam}
\end{figure}

Los mapas resaltan regiones que favorecen el género predicho según los
gradientes de la última etapa convolucional \cite{selvaraju2017}. En el acierto
de Amanita aparece activación en el entorno del ejemplar. En Auricularia, parte
de la señal coincide con los bordes de los cuerpos fructíferos. En el error de
Xylaria, el mapa resalta la zona donde se juntan el brazo y el tronco, una
posible señal ajena al ejemplar. La intensidad es relativa a cada foto.
El diagnóstico es cualitativo y corresponde a la rama visual compartida.

\section{Reproducción}
Las métricas por género de las seis variantes y los resultados ante
perturbaciones se incluyen en \texttt{reports/week\_11}. Para reproducir
la evaluación y compilar el informe desde la raíz del repositorio:

\begin{verbatim}
python scripts/train/evaluate_s11.py
python scripts/reports/build_s11_report.py
\end{verbatim}

\begin{thebibliography}{9}
\bibitem{selvaraju2017}
R. R. Selvaraju et al.
\href{https://openaccess.thecvf.com/content_iccv_2017/html/Selvaraju_Grad-CAM_Visual_Explanations_ICCV_2017_paper.html}{Grad-CAM: Visual Explanations from Deep Networks via Gradient-Based Localization}.
\emph{ICCV}, pp. 618--626, 2017.
\end{thebibliography}
\end{document}
